\chapter{Methodik}
Die wissenschaftliche Arbeit wird auf dem sechsstufigen Prozess des Design Science Research nach Pfeffers et al. \autocite{peffersDesignScienceResearch2007} basieren, welcher einen strukturierten Rahmen für die Entwicklung und Evaluation eines Artefakts bietet.

\section{Forschungsansatz: Design Science Research}
Im ersten Schritt wird hierbei das Problem identifiziert und die Motivation für die Arbeit beschrieben, wobei sich das Problem auf die fehlende Datenschutzkonformität aktueller Sprachassistenten bezieht und die Motivation der Arbeit in der Minderung gesammelter Daten für Nutzer liegt \autocite{peffersDesignScienceResearch2007}.

Anschließend werden genaue Anforderungen, also qualitative, aber auch quantitative Ziele, für eine Lösung unter Betracht der Problematik erstellt \autocite{peffersDesignScienceResearch2007}.

Mit diesen Grundbausteinen kann nun ein Artefakt, in diesem Falle ein Offline"=Sprachassistent, entwickelt werden \autocite{peffersDesignScienceResearch2007}.
Das Artefakt muss jedoch nicht unbedingt vollständig sein, aber mindestens die zuvor definierten Anforderungen erfüllen.

Dieses Produkt wird jetzt in verschiedenen Use Cases demonstriert \autocite{peffersDesignScienceResearch2007} (beispielsweise das Führen einer einfachen Unterhaltung auf leistungsstarker, aber auch ressourcenbeschränkter Hardware), bei dem ein oder mehrere der Probleme des ersten Schrittes gelöst werden.

Nah an die Demonstration angeknüpft wird das Artefakt dann gegen die Ursprungsproblematik evaluiert \autocite{peffersDesignScienceResearch2007}.  Im Zuge dessen werden die gemessenen Ergebnisse mit den anfänglichen Zielen, aber auch ähnlichen Lösungen verglichen. Im Falle des Sprachassistenten können beispielsweise die Antwortzeiten verschiedener Hardware gegen existierenden Cloud"=Lösungen verglichen werden.  Weiterhin werden im Evaluationsteil auch verschiedenste Use Cases für das Produkt vorgestellt und die Funktion des Sprachassistenten in diesen erläutert.

Abschließend werden im sechsten und letzten Schritt die Ergebnisse der Arbeit, sowie die Bedeutung des entwickelten Artefakts für die Forschung und Praxis kommuniziert. Dies beinhaltet unter anderem eine kritische Diskussion über die Limitationen, wie die begrenzte Rechenleistung bestimmter Hardware, sowohl auch der Herausforderungen, wie Modellaktualisierungen. Die gewonnenen Ergebnisse sollen letztendlich für professionelle und akademische, aber auch andere relevante Gruppen zugänglich gemacht werden. \autocite{peffersDesignScienceResearch2007}


\section{Evaluationskonzept}


\section{Begründung der Methodenauswahl}
